\documentclass[UTF8,12pt,a4paper]{ctexart}

\usepackage{amsmath}
\usepackage{cases}
\usepackage{cite}
\usepackage{graphicx}
\usepackage{enumerate}
\usepackage{algorithm}
\usepackage{caption}   % \caption*需要
\usepackage{subcaption} % 子图布局
\usepackage[noend]{algpseudocode} %algorithmicx
\makeatletter
\renewcommand{\fnum@algorithm}{\fname@algorithm}
\makeatother
\renewcommand{\algorithmicrequire}{\textbf{Input:}}
\renewcommand{\algorithmicensure}{\textbf{Output:}}
\usepackage[margin=1in]{geometry}
\geometry{a4paper}
\usepackage{fancyhdr}
\pagestyle{fancy}
\fancyhf{}
\usepackage{xcolor}

% ------------------- 自定义代码展示风格 -----------------
\usepackage{listings}
\usepackage{xcolor}
\lstdefinestyle{hspice}{
    language=C,
    basicstyle=\ttfamily\small,
    keywordstyle=\color{blue}\bfseries,
    commentstyle=\color{green!50!black},
    stringstyle=\color{brown},
    numbers=left,
    numberstyle=\tiny\color{gray},
    stepnumber=1,
    numbersep=5pt,
    backgroundcolor=\color{white},
    showspaces=false,
    showstringspaces=false,
    showtabs=false,
    frame=single,
    rulecolor=\color{black},
    tabsize=2,
    captionpos=b,
    breaklines=true,
    breakatwhitespace=true,
    escapeinside={\%*}{*)},
    morekeywords={param, include, temp, option, probe, DC, dc, measure, alter, VDD, GND, gnd, NMOS, PMOS,FIND, WHEN, sweep, abs, I, SUPPLY, lg, fin_height, fin_width, tran, ac, op, model, end, subckt,ends,global, nmos, pmos, inv, data, +}, % 根据需要添加更多HSPICE关键字
    comment=[l]{*} % 在这里设置注释
}

\lstset{style=hspice} % 只应用样式



% ------------------- 设置超链接，便于跳转 -----------------
\usepackage{enumitem}
\usepackage{hyperref}
\hypersetup{
    pdfborder={0 0 0},    % 消除超链接边框
    colorlinks=true,       % 将边框改为颜色标记（可选）
    linkcolor=black,       % 内部链接颜色（目录、引用等）
    citecolor=black,       % 引用颜色
    urlcolor=blue          % URL链接颜色（保持蓝色以便区分）
}




% ------------------------ 标题区 ----------------------------
\title{DIC Lab3 Report}
\author{
	姓名：\underline{冯峻}~~~~~~
	学号：\underline{523031910148}~~~~~~}
\date{\today}
\pagenumbering{arabic}

\begin{document}



% ------------------------ 页眉和页脚 ----------------------------
\fancyhead[L]{冯峻}
\fancyhead[C]{Lab3 Report}
\fancyfoot[C]{\thepage}

\maketitle
%\tableofcontents
%\newpage



% % -------------------------- 双图模板 ---------------------------
% \begin{figure}[htbp]
%     \centering
%     % 第一张子图
%     \begin{subfigure}[b]{0.45\textwidth} % 子图宽度为总宽度的45%
%         \centering
%         \includegraphics[width=\textwidth]{figure/wv_Task1_Fin_NFET.png} 
%         \caption{NMOS: $I_{\text{DS}}-V_{\text{DS}}$}
%         \label{fig:NFET_I_V}
%     \end{subfigure}
%     \hfill % 添加水平间距
%     % 第二张子图
%     \begin{subfigure}[b]{0.45\textwidth}
%         \centering
%         \includegraphics[width=\textwidth]{figure/wv_Task1_Fin_PFET.png} 
%         \caption{PMOS: $I_{\text{DS}}-V_{\text{DS}}$}
%         \label{fig:PFET_I_V}
%     \end{subfigure}
%     % 总图标题
%     \caption{$I_{\text{DS}}-V_{\text{DS}}$ curve}
%     \label{fig:FET_IV}
% \end{figure}


\end{document}
